\section{Conclusions}\label{sec:concl}

The exclusion region of Fig.~\ref{fig:map} is the result of this paper. For an exothermic pseudo-Dirac dark sector whose quark couplings follow the quark masses, a fixed xenon recoil yield and the Planck abundance determine the dark mass and coupling at every $(M_V,g_t)$, and the CMS all-hadronic $t\bar t$+MET search then excludes most of the resulting plane. Of the 470 admissible points, 381 are excluded at 95\% CL and 89 survive, 43 of them in the $M_V=20$~GeV column. At and above 100~GeV the entire scan is excluded, with expected limits reaching $\mu_{95}=0.00315$ at 200~GeV. Between 50 and 90~GeV the exclusion covers the whole scanned coupling range above $g_t\approx0.17$ and is bounded from below only; at 30 and 40~GeV a second boundary closes it at $g_t\approx2.7$--$3.0$; at 20~GeV the nominal exclusion is confined to an interval that the present Monte Carlo statistics cannot resolve.

The boundary has a simple physical reading. On the low-coupling side the signal becomes production-limited and the expected limit climbs through unity. On the high-coupling side the mediator decays increasingly into quarks, the invisible branching ratio at the largest scanned coupling falls from $0.89$ at 200~GeV to $2\times10^{-5}$ at 20~GeV, and the expected limit flattens and turns up: at 30 and 40~GeV this closes the window, from 50 to 90~GeV it turns up without reaching unity inside $g_t\le\sqrt{4\pi}$, and at 20~GeV the closure is not resolved. Where no closure occurs the excluded region runs to the largest couplings allowed by the narrow-width condition, and the production threshold alone carves out the surviving wedge at low mediator mass.

Three tests of different depth support this map. The relic equation was solved at both edges of the Planck window for all 470 admissible nodes, which leaves every one of them admissible and moves the dark mass by $0.07\%$ at the median; the only node the window can move across the boundary is one that was already unresolved. For the same nodes the signal was rescaled by the change in the invisible branching ratio, which shifts the limit by $0.02\%$ at the median. Dedicated endpoint samples at four nodes show that the efficiency carries no systematic shift beyond its statistical scatter, and that the boundary cases move by less than that scatter, one of them (139) sitting too close to unity for the present statistics to classify. Three extensions would sharpen the result: a full off-shell treatment extending the grid below $M_V=20$~GeV, a neural-network binning matched to the CMS extraction, and dedicated generation at the window edges for all surviving points, replacing the argument from measured shifts by a per-point verification. The central result stands on the present computation: within the scanned range the plane is excluded at and above 100~GeV, and the surviving wedge is a small and physically identifiable corner of the parameter space.
